\documentclass[../main.tex]{subfiles}
\begin{document}
%==========================================TIÊU ĐỀ TẬP========================================
\begin{table}[h]
  \begin{adjustwidth}{-1cm}{}
    \begin{tabular}{>{\centering\arraybackslash}p{6cm}|>{\raggedright\arraybackslash}p{12.5cm}}
      \multirow{5}{6cm}
      % Cột bên trái: ÉPISODE và số tập
      {\\[-40pt]
        \flaregothic\fontsize{20pt}{20pt}\selectfont \centering
        \color{eptype}ÉPISODE\color{black}\\
        \burbank\fontsize{72pt}{72pt}\selectfont
        \color{epnum}61\color{black}\\[6pt]
        \cafeteria\fontsize{20pt}{20pt}\selectfont\color{epnum}(5-5)\color{black}
        \\[18pt]
        \flaregothic\fontsize{14pt}{14pt}\selectfont
        \color{epattr}BIENVENUE, \uline{\color{Mtfk} Coco} !
        \color{black}
        \\}
       &
      % Dòng thứ nhất: tên tập tiếng Nhật
      {\fotskip\fontsize{18pt}{18pt}\selectfont \ruby{全}{ぜん}\ruby{員}{いん}ぶっ\ruby{潰}{つぶ}せ！
      } \\[6pt]
      % Dòng thứ hai: tên tập tiếng Anh
       & {\cafeteria\fontsize{18pt}{18pt}\selectfont Annihilate 'em All!} \\[4pt]
      % Dòng thứ ba: tên tập tiếng Việt
       & {\vnmsans\fontsize{18pt}{18pt}\selectfont Băng đảng} \\[8pt]
      % Dòng thứ tư: Nhân vật xuất hiện (theo đúng thứ tự debut: OC → Fubuki(17) → Marine(32) → Coco(34))
       & {\caviar{\fontsize{14pt}{14pt}\selectfont Nhân vật: \emp{kuon1} \emp{fubuki} \emp{marine} \emp{coco}}}
      \\[8pt]
      % Dòng cuối cùng: Ngày phát hành
       & {\continuum\fontsize{16pt}{16pt}\selectfont Ngày phát hành: 05/07/2020} \\[18pt]
    \end{tabular}
  \end{adjustwidth}
\end{table}
%=========================================TRUYỆN CHÍNH========================================
\color{black}

% Hộp tóm tắt
\txtbx[2]{{\color{vol} \uline{Tóm tắt:}} Kiryu Coco - nàng Rồng 3500 tuổi, thích đột kích, chửi thề bằng tiếng Anh và cưỡi khủng long - chính thức gia nhập \hll. Ngày đầu tiên của cô tại văn phòng đã trở thành một cuộc hỗn loạn không ngừng: từ những màn ``đột kích" lịch sự, đeo mặt nạ Jason, đến vác súng tên lửa và cuối cùng là phá tan ba ô cửa kính bằng một con T-Rex. Fubuki và Marine tham gia nhiệt tình, còn Kuon chỉ biết đứng nhìn và đếm thiệt hại. A-san xuất hiện, mắng một trận tơi bời, và cả ba bị trừ lương.}

\begin{center}
  \uline{(Người viết tập này có sử dụng bản dịch từ \href{https://www.youtube.com/@TomangclipHololivevedich}{Tớ mang clip \hll\ về dịch - TMCHVD}.)}
\end{center}

\pov{Hikawa Kuon}

Ôi trời. Năm tuần tiếp theo đây sẽ là năm tuần đầy mệt mỏi đây. Tôi đã sống sót qua đủ loại tình huống oái oăm với Thế hệ Thứ ba - những vụ cá khô bốc mùi, việc Marine tha rắn về, buổi chuẩn bị Comiket mùa đông hỗn mang với Pekora, và quảng cáo Othello có phần ``ảo ma'' của Noel-Rushia. 

Nhưng giờ đây, tôi sắp phải đối mặt với năm thành viên của \textit{holoForce} - hay còn gọi là Thế hệ Thứ tư, mỗi người một kiểu ``gây đau đầu'' khác nhau. Tôi thực sự nên được tăng lương. Ít nhất thì cũng nên được thưởng thêm một ngày nghỉ phép.

Tuần đầu tiên, chúng ta có gì nào?

Một cô nàng Rồng trẻ (ít nhất là trong giới Rồng, vài ngàn tuổi chỉ là mới lớn thôi) với tính cách nửa trẻ con, nửa quái kiệt. Trẻ con vì em ấy thích nghịch ngợm, cợt nhả, và dường như có một ``nút bấm'' khiến cô bật ra những câu chửi thề bằng tiếng Anh bất cứ lúc nào. Nhưng cũng có chút giang hồ vì cứ mở miệng ra là xả tiếng Anh như một rapper bị kích động. Tôi hiểu tiếng Anh, vì đã từng học qua vài năm ở trường, nên tôi biết rõ em ấy đang nói gì, và tôi đã chọn cách\dots\ không dịch lại cho những người không hiểu.

Một cô gái tới từ một thế giới rồng nào đó, nhưng tôi hay gọi đùa là ``USDA'' - United States of Dragon America. Nói sao nhỉ, em ấy có tất cả những đặc điểm của một công dân xứ Cờ Hoa của Long: giọng nói mạnh mẽ, phong thái tự tin, và đặc biệt là một nỗi ám ảnh kỳ lạ với văn hóa loài người. Tôi thề là em ấy có thể nói về ``kechi'' (tiết kiệm tiền), ``supacha'' (superchat) và ``peko'' suốt cả ngày mà không biết chán. Có lần tôi nghe em ấy nói chuyện với Fubuki về các loại hình thức tiết kiệm, tôi tưởng mình đang nghe một chuyên gia tài chính chứ không phải một Rồng.

Và, bởi vì là Rồng, em ấy cực kỳ tự hào về chủng tộc của mình. Dù vẫn giữ danh dự của một trong những sinh vật oai nghiêm nhất, nhưng nhờ khả năng kiểm soát nội lực mà em ấy có thể giữ hình dạng con người trong thời gian dài - vốn tiện lợi hơn rất nhiều so với việc đi lại trong văn phòng với một đôi cánh khổng lồ, nếu không muốn nói là nguyên một cơ thể cao bằng bảy, tám tầng. Tôi không muốn phải sửa trần nhà mỗi tuần đâu.

\chrint

Đúng, tôi đang nhắc đến thành viên đầu tiên của \textit{holoForce}, và cũng là nhóm trưởng, Kiryu Coco (\ruby{桐}{きり}\ruby{生}{う}ココ, \ruby{Ki-ri}{Đồng}-\ruby{ư}{Sinh}\ Cô-cô).

Coco là một nàng Rồng rất thẳng thắn, độc miệng, và thường khá om sòm - cái kiểu om sòm mà bạn có thể nghe thấy từ hành lang dù cánh cửa đã đóng chặt. Cô có khiếu hài hước thô lỗ và hay chửi thề (bằng tiếng Anh). Em ấy thích trêu chọc các thành viên khác trong các buổi phát trực tiếp của mình, chế giễu những thất bại của họ trong chương trình buổi sáng do em ấy dẫn, hoặc chế giễu trong các buổi hợp tác. Nói chung, nếu bạn là một Rồng, có lẽ bạn có cái quyền đó.

Em là một trong những VTuber đầu tiên của \hll\ - nếu không muốn nói là giới YouTuber ảo nói chung - phá vỡ ranh giới những gì VTuber có thể làm trên YouTube. Ví dụ như việc\dots\ phát trực tuyến về triệt lông hàng tháng (tôi vẫn không hiểu tại sao Rồng cần triệt lông), thảo luận về việc sử dụng kỹ thuật chạy nhanh trong nhà vệ sinh (RTA - Real Time Attack), hoặc kỹ thuật tấn công thời gian thực, thumbnail lật con chim (tôi xin miêu tả chi tiết), ám chỉ đến việc sử dụng ma túy (dù tôi tin đó chỉ là đùa), và dạy cho khán giả chủ yếu là người Nhật của mình cách sử dụng khác nhau của các từ chửi thề và tục tĩu bằng tiếng Anh. Đúng là một nhà giáo dục thế hệ mới.

\model{24}{l}{8cm}{coco}

Tuy nhiên, khi tình huống có chút chững lại, em ấy có thể rất nghiêm túc và trưởng thành. Coco rất từng trải, và thỉnh thoảng chia sẻ lời khuyên về các chủ đề nghiêm túc khác nhau với người xem và thậm chí là đồng nghiệp của em ấy khi được nhắc đến. Tôi đã thấy em ấy dỗ dành một đồng nghiệp đang khóc chỉ trong vài phút, rồi ngay sau đó chuyển sang chửi thề một streamer khác vì đã thua trong game. Đó là sự đa chiều của Rồng.

Như đã nói ở trên, Coco đôi khi sử dụng tiếng Anh và tiếng Nhật thay thế cho nhau. Ví dụ như câu cửa miệng ``\textbf{Good morning, motherfuckers!}'' (Chào buổi sáng, lũ khốn nạn!) bằng tiếng Anh, rồi ngay lập tức chuyển sang ``\textbf{Okite, okite, okite ku-da-sai!}'' (Dậy đi, dậy đi, dậy ngay!!) bằng tiếng Nhật. Sự chuyển đổi ngôn ngữ này đôi khi làm tôi choáng váng, nhưng cũng rất\dots\ Coco.

Thông thường, Coco có mái tóc dài màu cam được tạo hình chữ A với một lọn tóc vàng và bím tóc được tết lệch sang trái; trên đỉnh đầu có một sợi tóc ngố - kiểu tóc ngố của Rồng trông ngộ nghĩnh hơn của con người nhiều. Đôi mắt của em có màu hồng-tím rực rỡ, và trên đỉnh đầu có cặp sừng màu đỏ tím, trong đó có buộc nơ kẻ đen-trắng ở sừng phải, như thể em ấy đang trang trí cho một bữa tiệc của Rồng.

Trang phục của em ấy gồm một chiếc áo sơ mi có cổ khoét hở khe ngực, với cặp ngực được giữ cố định bằng một lớp lót - tôi không biết đó là thiết kế thời trang hay một chiếc áo chống đạn. Em ấy khoác một chiếc áo choàng đen có cổ tay áo màu đỏ tía và vạt mở ở ngực, một chiếc đầm váy ngắn màu kéo dài xuống gấu váy, được giữ cố định bằng nút thắt. Em đi một đôi tất bất đối xứng (một cao, một thấp), và đôi giày da mà học sinh hay đi. Em có đeo một chiếc vòng cổ rồng lớn có một viên ngọc đỏ, trông giống như một chiếc trâm cài lớn do vị trí của nó. Em ấy cũng đeo một chiếc bờm được trang trí bằng những ngôi sao bốn cánh lấp lánh, như thể vừa mới từ một buổi hòa nhạc của Rồng về.

\chrint

Mà, vấn đề lớn nhất của tôi là gì?

Là cô ấy lên văn phòng từ hôm nay.

Và tôi biết chắc chắn một điều: sự bình yên đã chấm dứt. Không phải là bình yên đã tồn tại ở đây trước đây, nhưng bây giờ nó còn chấm dứt nhanh hơn. Cơ mà văn phòng \hll\ này có bao giờ được yên ổn kể từ khi tôi bước chân vào đây đâu\dots\ Có lẽ tôi nên chuẩn bị sẵn thuốc đau đầu, và cà phê, và thêm một bộ giáp chống chửi thề.

\il

Một ngày đầu tháng Bảy. Nắng nóng như đổ lửa, tiếng ve kêu râm ran, và tôi đã có một dự cảm chẳng lành từ lúc bước chân vào văn phòng.

Tôi không biết số phận đã sắp đặt thế nào, nhưng bình yên thực sự đã chấm dứt từ giây phút cô nàng Rồng mới nhập hội đặt chân vào văn phòng. Một cảm giác như có một cơn bão đang tụ lại phía chân trời, và tôi là người duy nhất nhận ra điều đó.

Ngày đầu tiên làm việc của em ấy trùng với một mùa hè oi bức - cái nóng đến nỗi quạt chạy hết công suất cũng chỉ đủ để không bị\dots\ tan chảy. Tôi vừa bước vào văn phòng, chưa kịp nhấp ngụm cà phê đầu tiên, thì đã nghe thấy một câu gây sốc phát ra từ phía bàn làm việc mới: ``Tự nhiên em muốn làm một cuộc đột kích ghê!''

Tôi dừng tay trên bàn phím, chậm rãi ngẩng đầu lên và nhìn về phía nguồn phát ra câu nói. Và tất nhiên, không ai khác chính là cô nàng Rồng - Coco. Em ấy đang đứng giữa phòng, hai tay chống nạnh, đôi mắt sáng rực như thể vừa phát minh ra một trò chơi mới.

Bên cạnh em ấy, Marine - người luôn có khẩu vị với sự hỗn loạn ngoài sự dâm đãng vốn có - cười khoái chí, mặt như đứa trẻ vừa thấy một quả bom mới. Cô tiếp lời ngay, giọng đầy phấn khích: ``Ái chà, em cũng khùng quá nhỉ. Thích thì cứ triển thôi. Có khi làm nên chuyện đấy.''

Nhưng vừa dứt câu, tôi thấy Coco tặc lưỡi. Một cái tặc lưỡi nhỏ, chỉ là một âm thanh nhẹ, nhưng đủ để làm thay đổi bầu không khí trong phòng. Marine nhướng mày, nhìn thẳng vào cô nàng hậu bối với ánh mắt sắc như dao: ``Hả? Em vừa mới tặc lưỡi với chị đấy à?''

Không khí bỗng trở nên căng thẳng như dây đàn. Tôi thở dài. Lại nữa rồi. Chưa đầy một phút từ khi đến văn phòng, cô Rồng này đã tính gây chiến với Senchou rồi sao? Tôi nhấp một ngụm cà phê, chuẩn bị tinh thần cho một trận khẩu chiến.

Và đúng lúc tôi tưởng hai người sẽ nổ ra một trận khẩu chiến nảy lửa - thứ mà tôi đã chứng kiến quá nhiều lần - thì Fubuki xuất hiện. Cô từ đâu bước ra (có lẽ từ sau cánh cửa tủ?), trên môi nở nụ cười tinh quái, giọng đầy hào hứng: ``Vậy bọn mình sẽ là `Băng đảng \hll' nha\~{} Chị sẽ là kẻ địch!''

\img{5-5a}

Ngay từ giây phút đó, tôi nhận ra một cơn bão sắp ập đến. Lần nữa. Làn sóng hỗn loạn mới đang hình thành, và tôi chỉ có thể đứng nhìn mà chẳng thể ngăn cản được.

Mắt của cô Rồng trẻ sáng lên như thể vừa được bật chế độ chiến đấu - những tia lửa lóe lên trong đôi mắt hồng-tím của em. Tôi có thể cảm nhận nguồn năng lượng nguy hiểm lan tỏa khắp căn phòng, như một cơn gió nóng trước bão.

Tôi nhanh chóng đặt tách cà phê xuống, tay tiếp tục gõ máy tính viết báo cáo, nhưng giọng thì đầy cảnh cáo, với hy vọng ngăn được thảm họa: ``Đừng có mà phá văn phòng lên nha, mấy đứa! Anh mới thay cửa sổ tuần trước.''

Ba cô gái đồng thanh đáp lại, kéo dài giọng đầy tinh quái, như ba nàng tiên của sự hỗn loạn: ``Vâng\~{}''

Tôi không biết có nên tin vào chữ ``vâng'' đó không. Mà không, chắc chắn là chẳng nên tin rồi. Cái giọng ``vâng'' kiểu đó thường có nghĩa là ``chúng em sẽ làm theo cách của chúng em, nhưng sẽ cố gắng không để anh phát hiện.''

\ct

Lần ``đột kích'' đầu tiên của Coco diễn ra theo một cách mà không ai lường trước được - kể cả tôi, người đã nghĩ rằng mình đã chuẩn bị cho mọi tình huống.

Cánh cửa mở ra một cách\dots\ lịch sự và nhẹ nhàng. Không đạp cửa, không hét lớn, không có hiệu ứng âm thanh đặc biệt. Chỉ là một cánh cửa mở ra từ từ, như một người đưa thư bước vào.

Và rồi, giọng nói nhỏ nhẹ, gần như thì thầm của Coco vang lên: ``Đột kích đây\dots''

Tôi còn chưa kịp phản ứng thì đã nghe một tiếng phê phán thẳng mặt từ Marine, giọng đầy thất vọng: ``Đang vào vai người đưa thư hả má trẻ!? Có ai đột kích kiểu đó không!?''

Tôi bật cười đến mức phải ngừng gõ máy tính. Thế quái nào mà lại có kiểu đột kích như thế này!? Ngay cả trong mơ tôi cũng không nghĩ đến. Tôi đã tưởng tượng ra cảnh cửa bị đạp tung, tiếng hét vang trời, và có lẽ là một vài mảnh gỗ bay khắp phòng.

``Senchou nói đúng đấy,'' tôi nói, vẫn chưa hết cười. ``Đột kích kiểu thế thì địch giết em con mẹ hàng lươn rồi. Có khi chỉ kịp hỏi `có thư gì không?' là bị xiên luôn.''

Marine còn tận tình chỉ dạy, vẽ ra một ví dụ thực tế bằng cách đứng dậy, vung chân lên trời như một võ sĩ karate đang tập: ``Ít nhất cũng phải đạp cửa chứ! Một cú đạp mạnh, cửa bật tung, người nhảy vào, hét lớn, thế mới gọi là đột kích!''

\img{5-5c}

``Hoặc dùng tay đập cửa cũng được,'' tôi gật đầu tán thành ``Nhưng mà đừng làm ở đây, hỏng cửa đấy. Anh không muốn phải báo cáo lên trên là `cửa phòng bị Rồng đạp vỡ trong một buổi đột kích tập thể' đâu.''

Nhưng nàng Rồng kia, thay vì tiếp thu, lại núp sau cánh cửa, lắc đầu nguầy nguậy, nở nụ cười vô tư như thể vừa nói một câu châm ngôn sống: ``Nói không với bạo lực anh chị à. Đột kích cũng có văn hóa của nó.''

Tôi và Marine quay sang nhìn nhau. Cả hai chúng tôi đồng thanh bật lại ngay lập tức, giọng đầy bất lực: ``Thế mà cũng bày làm đột kích!?''
``Xem cái đứa muốn chơi đột kích nói gì kìa!''

Coco câm nín trong giây lát, hai mắt chớp chớp, rồi chỉ lè lưỡi một cách nghịch ngợm - một biểu cảm vừa đáng yêu vừa thách thức. Đúng là cứng đầu hết thuốc chữa. Cô Rồng này có thể bẻ cong lý luận như bẻ cong một thanh sắt.

Nàng Thuyền trưởng khoanh tay, tiếp tục mắng với giọng của một người đã từng làm đủ mọi trò điên rồ: ``Liệu mà làm cho đàng hoàng vào! Lần sau đạp cửa, hét lớn, và nhảy vào, chứ đừng mở cửa từ từ như kẻ trộm!''

Và thế là, nàng Rồng lại đóng cửa, chuẩn bị cho lượt ``đột kích'' tiếp theo. Tôi có thể nghe thấy tiếng em ấy thì thầm điều gì đó như ``lần này sẽ làm đúng\dots''

\il

Ngay lúc đó, Cáo trắng - xuất hiện từ đâu đó (tôi đoán là cô ấy luôn xuất hiện đúng lúc để bổ sung câu chốt), như thể đã chờ sẵn để góp một câu: ``Tại \hll, bọn mình về cơ bản là nhận được thư mỗi ngày.''

\img{5-5d}

Tôi đứng bên cạnh, thở dài, rồi bổ sung thêm một chi tiết quan trọng mà Fubuki đã cố ý quên, một sự thật phũ phàng mà cô ấy thường bỏ qua: ``Toàn hóa đơn tiền điện, tiền nước với phí sửa chữa hết đó. Đôi khi là thư đòi nợ nữa, nhưng chúng tôi không nói về điều đó.''

Cáo trắng cười khì, như thể đang nhớ lại những lá thư tình ngọt ngào, trong khi tôi thì bĩu môi, tỏ rõ vẻ không hài lòng. Không phải ai cũng có niềm vui nhận thư như em đâu, Fubu-chan à. Em thích đọc hóa đơn, nhưng anh thì không.

``Đấy, Coco-chan này,'' tôi nói, vừa chỉ tay vào cô Rồng đang chuẩn bị đột kích lần hai, ``nếu em muốn đột kích, hãy đột kích vào đống hóa đơn kia. Chúng đang chờ được giải quyết từ tuần trước.''

Coco ngó đầu ra khỏi cánh cửa, nhìn đống giấy tờ trên bàn tôi, rồi nhăn mặt: ``Chừng đó em thà đột kích phòng của Kanatan còn hơn.''

\il

\ct

Lần ``đột kích'' thứ hai của cô Rồng đã nâng cấp lên một tầm cao mới - và lần này, tôi thực sự giật mình. Không phải vì sợ, mà vì sự đầu tư công phu đến mức\dots\ đáng kinh ngạc.

Cửa bị đạp tung ra với một lực đủ mạnh để tôi nghe tiếng bản lề kêu răng rắc. Tôi có thể thề là cả tòa nhà cũng rung lên. Và rồi, từ phía sau cánh cửa, một bóng đen xuất hiện.

\img{5-5e}

Một kẻ mang mặt nạ Jason Voorhees, tay cầm cưa máy, xuất hiện. Cưa máy thậm chí còn phát ra tiếng động giống thật - tôi không biết em ấy đã lấy nó ở đâu, nhưng nếu đó là đồ chơi thì tôi phải công nhận nó đáng sợ thật.

Tôi giả vờ bật ngửa trên ghế, hét lớn như để hùa theo không khí: ``Ối giời ơi, sợ thế! Ai đó gọi cảnh sát đi!''

Cơ mà\dots\ Không ai bảo ai, cả tôi lẫn Marine đều cảm thấy có gì đó sai sai. Jason không bao giờ cầm cưa máy. Đó là Leatherface, một kẻ giết người khác. Sự thật lịch sử đang bị xuyên tạc ngay trước mắt tôi.

Thay vì giả vờ hoảng hốt như tôi, Thuyền trưởng phán một câu lạnh tanh, giọng như một chuyên gia phim kinh dị đang chỉ ra lỗi kỹ thuật: ``Khi nào tới thứ Sáu hẵng quay lại nha! Jason mà lại xuất hiện cuối tuần thì không hợp lý tí nào!''

Tôi hiểu ý của em ấy ngay lập tức, liền chêm thêm: ``Thứ Sáu ngày 13 ấy. Đó mới là ngày của hắn. Còn bây giờ là Chủ Nhật, không ai sợ ngày này cả.''

Nhưng nàng Rồng, thay vì tận hưởng vai diễn của mình, lại hổn hển tháo chiếc mặt nạ ra, vẻ mặt đầy bất mãn, mồ hôi lấm tấm trên trán: ``Cái này khó thở ghê vậy đó. Em tưởng mua loại thoáng khí, ai ngờ\dots''

Marine nghe vậy bóp trán, cáu kỉnh đáp, giọng đầy thất vọng: ``\textbf{Đã bảo là làm cho đàng hoàng vào mà!} Jason mà phải than thở vì khó thở thì thua cả xác chết!''

Tôi bật cười, lắc đầu. Đúng là một cô Rồng đầy bất ngờ.

\il

Ngay lúc đó, Cáo trắng lại xuất hiện - lần này với mặt nạ Jason vẫn còn trên tay, tỏa ra một bầu không khí tri thức như một giáo sư đang chuẩn bị lên lớp. Tôi thầm đếm ngược trong đầu.

Và đúng như dự đoán, em ấy bắt đầu chương trình giáo dục của mình, với giọng như đang đọc một bài giảng quan trọng: ``Cái ông trong phim Thứ Sáu ngày 13 ấy\dots''

\img{5-5f}

Không để em ấy làm màu quá lâu, tôi đáp nhanh để tiết kiệm thời gian, như một học sinh đã quá quen với câu chuyện: ``Gã Jason cho mấy ông chưa biết. Chúng ta đang nói về hắn.''

Fubu-chan gật gù rồi nói tiếp, vẫn giữ nguyên cái điệu giảng bài đầy tự hào: ``\dots\ chưa bao giờ cầm cưa vào tấn công chỗ này. Trong phim, hắn chỉ dùng dao, rìu, hoặc tay không thôi.''

Tôi khoanh tay, nhướng mày nhìn sang em ấy, rồi chốt hạ một câu đầy châm biếm: ``Hắn còn chưa xuất hiện bao giờ trong này, huống gì cầm cưa. Em muốn cả lũ lên bàn thờ ngắm mâm ngũ quả hả? Jason mà đột nhập vào văn phòng của mình thì chắc chỉ để đòi hóa đơn thôi.''

Fubuki nghe vậy cũng chỉ cười khì, gãi đầu: ``Đ-Đúng ha\dots\ Có lẽ em nên xem lại phim trước khi kể fun fact\dots''

Tôi lắc đầu, vừa cười vừa thầm nghĩ: ``\textit{Đúng là Fubu, lúc nào cũng muốn khoe kiến thức, nhưng mà kiến thức kiểu này thì thà không khoe còn hơn.}''

\il

\ct

Lần ``đột kích'' thứ ba, Coco tiếp tục thử nghiệm với một phong cách hoàn toàn mới. Có vẻ như em ấy đang luyện tập, và mỗi lần đều có sự tiến bộ rõ rệt.

Cửa bật tung ra với một lực mạnh hơn, bản lề kêu răng rắc lần nữa. Lần này, không chỉ đạp cửa, nàng Rồng còn vác theo một khẩu súng phóng tên lửa - tất nhiên là đồ giả, hoặc ít nhất là tôi nghĩ vậy, nhưng nhìn sơ qua thì cũng đáng sợ thật. Nó có đèn LED nhấp nháy và phát ra tiếng bíp, như một món đồ chơi công nghệ cao.

\img{5-5g}

Vẫn giữ nguyên cái giọng vui tươi và vô tư đó, em ấy đưa khẩu súng lên nhắm về phía tôi và tuyên bố một câu đầy kịch tính: ``Nhớ lúc ta bảo là sẽ giết ngươi sau cùng chứ? Giờ thì ngươi sẽ phải trả giá!''

Nhưng rồi, đột nhiên giọng em ấy trầm hẳn đi, mang theo một sắc thái nguy hiểm khiến tôi suýt giật mình, nếu tôi không biết rằng em ấy chỉ đang diễn: ``Xạo chó đấy. Tôi sẽ giết ngươi trước.''

Tôi đứng chắp tay sau lưng, gật gù đánh giá như một vị giám khảo khó tính: ``Hơi sai thoại một tí, nhưng thế là OK rồi. Lần sau nhớ đọc kịch bản trước khi diễn.''

Lúc này, Marine và Fubuki cũng bắt đầu nhập vai, khiến cả phòng trở thành một sân khấu ứng tác đầy bất ngờ.

Thuyền trưởng giơ bộ đàm - cũng không chắc là hàng xịn hay hàng đểu, nhưng trông có vẻ rất chuyên nghiệp - giả vờ đang liên lạc với sở chỉ huy. Cô bình thản thông báo, giọng như một cảnh sát trưởng trong phim hành động: ``Báo cáo, kẻ xâm nhập đã được xác định. Hắn ta là một tên dâm tặc đô con, vũ trang hạng nặng, hết!''

\img{5-5h}

Không để Senchou độc diễn, tôi cũng vào vai, nhấn mạnh giọng trầm đầy quyền uy, như thể đang chỉ huy một đội quân tinh nhuệ: ``Chuẩn bị sẵn sàng tác chiến, phương án B. Triển khai bẫy bánh pudding nếu cần, hết!''

Nhưng Coco thì vẫn giữ vẻ hồn nhiên đến khó tin, chẳng hề nao núng trước sự hợp tác hoàn hảo của tôi và Marine. Em ấy hạ súng xuống, nhìn chúng tôi với vẻ ngây thơ: ``Thật ra thì em chưa xem bộ phim đó nữa. Em chỉ xem mấy clip cắt trên YouTube thôi.''

Senchou lập tức phán một câu cụt lủn vào bộ đàm, giọng đầy thất vọng: ``Làm cho nghiêm túc đi. Hết.''

\il

Ngay lúc đó, Fubuki - với khuôn mặt lấm lem như quân nhân ngoài chiến trường (cô bôi bột mì lên mặt từ lúc nào không rõ) - lại xuất hiện, mang theo một ``fun fact'' mới để khoe sự hiểu biết của mình.

\img{5-5i}

``Trong bản gốc (phim Âu Mỹ), thì khi kẻ xấu bảo: `Nếu muốn chuộc lại đứa con thì ngươi phải hợp tác', đúng chứ? Nhưng thay vì nói `Không', trong bản lồng tiếng Nhật thì lại nói: `OK!' đó.''

Tôi liếc nhìn nàng Rồng, và em ấy cũng liếc lại tôi. Cả hai cùng cau mày đầy nghi ngờ, như thể vừa nghe một tin tức không thể tin nổi.

``Thật hả?'' tôi hỏi, giọng đầy hoài nghi.

``Có vụ này luôn sao?'' Coco đồng thanh, với vẻ mặt cũng không kém phần hoang mang.

Không ai nói gì thêm một lúc. Bầu không khí im lặng bao trùm, chỉ có tiếng quạt chạy và tiếng ve ngoài cửa sổ.

Rồi tôi lắc đầu, cắt ngang luôn với giọng đầy quyết đoán: ``Không, không. Sai rồi. Fubu-chan, em đã bị ai lừa đấy.''

Coco gật gù, tiếp lời ngay tức khắc: ``Đúng á. Mấy bản lồng tiếng Nhật ấy, có thể thay đổi thoại để phù hợp với khẩu hình, nhưng mà đổi hẳn `Không' thành `OK' là lạ lắm. Nó sẽ làm hỏng toàn bộ mạch cảm xúc của phim đó, Fubuki-paisen à.''

Tôi khoanh tay, suy ngẫm một chút trước khi tiếp tục: ``Có thể là có trường hợp này thật, nhưng chắc chắn không phải một câu thoại mang tính căng thẳng như vậy. Lồng tiếng Nhật hay thay đổi nội dung để khớp với văn hóa hơn, nhưng không phải theo kiểu đó. Nó sẽ làm phim mất đi tính kịch tính.''

Fubu cười khì, giơ tay lên như đang vẫy cờ trắng, vẻ mặt đầy ngượng ngùng: ``Thì em chỉ nói `\eng{fun fact}' thôi mà\~{} Không cần phải soi mói đến thế đâu.''

Tôi thở dài, lườm nàng Cáo trắng bằng ánh mắt trách móc, giọng đầy bất lực: ``\eng{Fun fact} mà sai thì nó chẳng còn \eng{fun} nữa, mà nó thành \eng{wrong fact} cmn luôn rồi! Em muốn làm giáo viên thì nhớ kiểm tra lại thông tin trước khi dạy người ta chứ!''

Coco cười phá lên, giọng vang khắp phòng, còn Fubuki thì chỉ nhún vai một cách vô tội, như thể chẳng có gì sai.

\il

\ct

Lần ``đột kích'' thứ tư, nàng Rồng chơi lớn. Và tôi không nói đùa. Cô ấy đã cưỡi một con khủng long bạo chúa - một con T-Rex to tướng, xanh lè, đủ cao để đâm thủng trần nhà nếu nó nhảy lên - xông thẳng vào văn phòng.

\img{5-5j}

Em ấy reo lên đầy phấn khích, giơ hai tay lên trời như một nữ chiến binh thời tiền sử: ``\textbf{Công viên kỷ Jura! Mấy người chưa chuẩn bị cho quả này đâu nhỉ? Hãy chào đón thời đại mới của các cuộc đột kích!}''

Cảnh tượng trước mắt khiến hai con người tự nhận là ``nghiêm túc nhất phòng'' (tôi và Marine) đổ mồ hôi hột. Tôi mắt chữ A mồm chữ O, hoảng hốt thốt lên: ``Em kiếm đâu ra con khủng long này vậy!? Nó có vẻ là đồ chơi, nhưng mà đồ chơi to cỡ này thì cũng là hàng hiếm!''

Senchou cũng không khá hơn, giọng lạc cả đi, tay chỉ về phía con T-Rex: ``Mà ai đâu biết được là cưng vác nguyên thứ đó tới chứ!? Nó giống thật quá, tưởng em đi cướp bảo tàng về!''

Trái ngược với sự ngỡ ngàng của chúng tôi, Fubu - như mọi khi - lại là người nhiệt tình nhất. Cô nhảy tót lên trước đầu con khủng long, mắt sáng rực, giọng như đứa trẻ vừa thấy đồ chơi mới: ``\textbf{Đù! Hàng thật luôn à!? Nó có chạy bằng pin không? Hay em dùng cưỡi thật?}''

\img{5-5k}

Sự phấn khích quá đáng của Fubuki khiến Senchou trợn mắt quay sang nhìn tôi, rồi hỏi đầy bất lực, như thể đang cầu cứu: ``Không có fun fact nào nữa à!? Tôi cần một thứ gì đó để giải thích hiện tượng quái dị này!''

Tôi cười nhẹ, giơ tay chỉ về phía cô nàng đang cưỡi trên con khủng long, giọng đầy vẻ ``đến lượt em'': ``Có chứ. Mời Coco-chan. Em ấy có vẻ đã chuẩn bị một bài thuyết trình đấy.''

Nàng Rồng khoanh tay lại, gật gù với vẻ hiểu biết, như một giáo sư chuẩn bị lên lớp, rồi giải thích đầy tự hào, giọng vang khắp phòng: ``Mục tiêu của đột kích chính là để tiêu diệt băng đối thủ. Đó là quy luật cơ bản của mọi cuộc đột kích. Nếu không có mục tiêu thì không gọi là đột kích, mà là\dots\ đi dạo.''

Lời giải thích của em ấy, đi kèm với tiếng gầm của con khủng long bạo chúa (có lẽ là loa gắn trong bụng) khiến Fubuki giật thót, dựng cả đuôi lên. Còn tôi thì chỉ nhún vai, bình thản hỏi: ``Nghe hiển nhiên quá. Thế mục tiêu là gì? Ai là đối thủ? Hay chỉ là muốn có cớ để phá hoại?''

Coco cười toe toét, đôi mắt hồng-tím sáng lên đầy tinh quái. Em ấy nhìn về những ô cửa sổ trước mặt, như thể đang nhìn vào một tương lai đầy bạo lực, rồi đáp lại đầy hồn nhiên: ``Phá cửa kính có gia huy của băng đối thủ trên đó! Chúng ta sẽ để lại dấu ấn của mình!''

Tôi mở to mắt, kinh ngạc thốt lên: ``\dots\ Hả? Cái gì cơ? Gia huy ở đâu ra? Văn phòng này làm gì có gia huy?''

Như bị kích thích, cả ba người con gái quay đầu lại nhìn về phía cửa kính. Tôi chỉ kịp nhăn mặt lại, thở dài khi nhận ra điều sắp xảy ra: ``Không phải lại nữa chứ\dots\ Mấy đứa đừng có\dots''

Nhưng đã quá muộn.

Trong tích tắc, cả ba cô gái cùng với con khủng long kia - bằng một cách thần kỳ nào đó - lao thẳng về phía ba ô cửa kính.

*RẦM!* *XOẢNG!* *CHOANG!* 
Ba tiếng vỡ vang lên, kinh thiên động địa. Kính vỡ tan tành, những mảnh vụn lấp lánh rơi xuống đường dưới ánh nắng chói chang.

\img{5-5l}

Họ cùng nhau hét lớn, giọng hòa làm một với tiếng gầm của khủng long: ``\textbf{Đột kích nào!!}''

Tôi lắc đầu, vò trán thở dài, mắt nhìn theo những mảnh kính đang rơi: ``Mấy đứa hiểu sai nghĩa `đột kích' rồi đó. Đột kích là bất ngờ xông vào, chứ không phải là phá cửa kính mỗi ngày. A-san mà thấy cái này thể nào cũng tức điên cho mà xem\dots''

Và ngay từ thời điểm đó, định mệnh của họ đã an bài.

Ngay khi tôi vừa dứt lời, một giọng hét đầy giận dữ vang lên từ phía sau, đủ sức làm rung cả những mảnh kính còn sót lại: ``\textbf{Mấy con dở hơi này!!}''

Tôi không cần quay lại cũng biết ai vừa xuất hiện.

A-san - với vẻ mặt đỏ bừng như cà chua chín, khắp người tỏa ra sát khí ngùn ngụt, đứng đó như một vị thần trừng phạt sắp giáng xuống cơn thịnh nộ. Trong tay cô là một tập hồ sơ dày cộp, mà tôi đoán chính là báo cáo thiệt hại mà cô đang viết dở.

Ba cô gái đang phấn khích bỗng cứng đờ như tượng đá, mặt tái mét, tay chân đông cứng. Con khủng long vẫn đứng đó, bất động, như một nhân chứng thầm lặng.

Còn tôi? Tôi chỉ chẹp miệng, khoanh tay lại, bình thản thở dài, như thể vừa dự đoán được một kết cục không thể tránh khỏi: ``\eng{Mẹ nó,} biết ngay mà. Lúc nào cũng vào đúng thời điểm.''

\begin{figure}[H]
  \centering
  \begin{subfigure}[t]{0.45\textwidth}
    \centering
    \includegraphics[width=8cm]{../../../../Image/Book5/5-5m1.png}
  \end{subfigure}
  \quad
  \begin{subfigure}[t]{0.45\textwidth}
    \centering
    \includegraphics[width=8cm]{../../../../Image/Book5/5-5m2.png}
  \end{subfigure}
\end{figure}

A-san chỉ tay về phía ba kẻ phá hoại, giọng như sấm, vang vọng khắp căn phòng: ``\textbf{Lũ bay tính đền kiểu gì đây hả!? Cửa sổ mới thay tuần trước! Vừa tuần trước! Mấy đứa biết một ô cửa sổ công ty giá bao nhiêu không!?}''

Tôi nhún vai, lắc đầu, ra vẻ bất lực trước độ ngáo của ba đứa kia, giọng đầy thương cảm: ``Thế này anh cũng không đỡ được mấy đứa rồi. Bán cả khủng long cũng không đủ trả tiền kính luôn dấy chứ.''

Cuối cùng, cả ba cô gái - Coco, Marine và Fubuki - đồng loạt cúi rạp người, giọng hòa làm một, đầy sám hối: ``\textbf{Chúng em xin lỗi!!}''

A-san vẫn đứng đó, hai tay chống nạnh, thở dốc, mắt vẫn trừng trừng nhìn ba kẻ phá hoại. Còn tôi, tôi chỉ biết lắc đầu, vừa cười vừa thở dài, nghĩ thầm: ``\textit{Chào mừng đến với \hll, nơi mà mỗi ngày đều là một cuộc đột kích mà chẳng một ai trong này biết trước được.}''

%==========================================TRUYỆN PHỤ=========================================
\omk

Sau trận đại náo cửa kính, ba kẻ phá hoại bị cấp trên của tôi - A-san - mắng cho một trận tơi bời. Tiếng mắng vang vọng khắp hành lang, đủ để làm những nhân viên khác trong tòa nhà phải co rúm lại.

A-san nhìn tôi chằm chằm, giọng đầy trách móc, tay vẫn cầm tập hồ sơ thiệt hại mà cô đã lập xong từ hôm qua: ``Hikawa-san, cơ sở vật chất bị hủy hoại quá nhiều rồi đấy. Cửa sổ tuần trước mới thay, giờ lại vỡ. Giờ cậu thì tính sao đây?''

Tôi bình thản trả lời, như thể đã có sự chuẩn bị từ trước - vì thực ra tôi đã lên kế hoạch cho tình huống này từ lúc Coco bước vào văn phòng: ``Yên tâm chị, em có ghi ai phá cái gì và đền thế nào rồi. Coco là người mới nên được tha lần này, nhưng Marine và Fubuki mỗi người sẽ bị trừ 15\% lương để đền tiền kính. Em cũng sẽ góp một phần, vì đã không ngăn cản kịp thời.''

Chị ấy nhướng mày, hơi bực bội, giọng vẫn chưa nguôi: ``Vậy còn chuyện cậu cứ để họ mặc sức tung hoành thì sao? Không có cách nào kiềm chế bọn họ à? Cậu là quản lý, phải có trách nhiệm chứ.''

Tôi nhún vai, thở dài một hơi, nhìn thẳng vào mắt chị ấy. Trong mắt tôi, một sự mệt mỏi chất chứa từ những lần chứng kiến các cô nàng \hll\ làm loạn: ``Việc này bọn họ phải tự làm thôi. Một mình em nhắc miệng không ăn thua. Thậm chí là phạt lương rồi, họ cũng có thèm nghe đâu. Em đã thử mọi cách: khuyên nhủ, đe dọa, phạt lương, thậm chí là mời họ đi ăn để mua chuộc - vẫn không hiệu quả.''

A-san chỉ biết đưa tay lên trán, thở dài trong ngao ngán, hai mắt nhắm nghiền như đang cầu nguyện: ``Thế thì phải làm sao? Cứ để họ phá hoại mãi sao?''

Tôi vỗ nhẹ vai chị ấy, cười nhẹ, giọng đầy vẻ chấp nhận số phận: ``Thôi nào, \hll\ là như thế, không loạn thì làm gì gọi là \hll! Nếu một ngày nào đó văn phòng yên tĩnh, em mới lo.''

Chị ấy thoáng ngạc nhiên với câu nói thốt ra từ một người như tôi - vốn dĩ coi sự nghiêm túc là số một - giờ đây cũng phải đối diện với sự thật nghiệt ngã này: các tài năng của \hll\ khi đến văn phòng này sẽ luôn luôn có chuyện. Thậm chí, tôi nghĩ rằng họ đến đây để phá hoại nhiều hơn là để làm việc.

Chị ấy nhìn tôi đầy bất lực, nhưng cũng đành gật đầu chấp nhận: ``Tôi nghĩ vậy\dots\ Chúng ta chắc phải tìm loại kính nào bền hơn nữa mới chịu.''

\ct

Hậu quả của trận ``đột kích'' phá cửa kính đó thực sự không thể nào chua chát hơn:

Marine và Fubuki, mỗi người bị trừ 15\% lương để đền tiền cửa kính. Con số này tương đương với một tháng ăn trưa của họ, nhưng ít nhất cũng đủ để A-san nguôi giận.

Coco thì đây là lần đầu tới văn phòng nên được tha bổng - theo đúng luật ``người mới được miễn trừ'' nhưng tôi đã cảnh báo em ấy rằng lần sau sẽ không có ân huệ đó.

Bản thân tôi cũng phải trích lương ra để góp vào, vì đã không ngăn cản kịp thời. Một phần tư tháng lương của tôi đi vào quỹ sửa chữa văn phòng.

Sau vụ đó, Thuyền trưởng và Cáo trắng tạm thời ngoan hơn. Marine thậm chí còn tự nguyện mang trà sữa cho cả phòng trong một tuần để ``chuộc tội''. Fubuki thì im lặng học thuộc các fun fact có thật để tránh bị tôi soi mói.

Nhưng nàng Rồng Coco? Không hề có dấu hiệu dừng lại. Dù bị mắng thế nào, cô nàng vẫn quậy phá như thường, khiến cả tôi và A-san phải đau đầu. Em ấy như một cơn bão không thể ngăn cản, và tôi bắt đầu hiểu rằng việc ``ngăn cản'' Coco cũng vô nghĩa như ngăn cản sóng biển.

Nhìn nàng Rồng vô tư quậy tung cả văn phòng - lần này là đang cố gắng cưỡi con khủng long qua cánh cửa vừa được sửa - A-san lẩm bẩm: ``Hy vọng là em ấy không phá tung cả cái văn phòng này\dots''

Tôi nhìn xung quanh, thở dài đáp lại, mắt nhìn theo Coco đang tập trung điều khiển con T-Rex: ``Chị nên nói là tất cả mọi người mới đúng đó. Không chỉ Coco-chan, mà Marine-chan, Fubuki-chan, và cả những người chưa đến nữa. Công ty này là một tập hợp những cơn bão.''

Chị ấy xoa cằm, suy tư: ``Chắc ta phải tìm cửa kính nào bền hơn nữa mới chịu. Kính cường lực, kính chống đạn, hay kính chịu được khủng long?''

Tôi gật đầu ngay lập tức, giọng đầy đồng tình: ``Đồng ý. Quả đúng là\dots''

Cả hai cùng ngẩng lên nhìn nhau, đồng thanh nói trong tiếng thở dài, như một câu thần chú bất thành văn: ``Ta đã triệu hồi một con báo rồi nhỉ\dots''

\il

Nhưng sóng gió vẫn chưa dừng lại ở đó. Đúng hơn, đây mới chỉ là mở đầu trong chuỗi năm tuần đón thành viên mới mà thôi. Một cơn bão vừa qua, nhưng còn bốn cơn bão khác đang hình thành.

Tôi cầm lấy tập hồ sơ của \textit{holoForce} (Thế hệ thứ Tư) trên mặt bàn mà A-san đã đưa cho tôi từ tuần trước để đọc sơ qua, từ đó chuẩn bị cho mọi kịch bản quậy phá mà họ có thể tạo ra. Bìa hồ sơ màu đỏ, bên trong là những trang giấy đầy thông tin về các thành viên.

Cái tên tiếp theo trong số bốn người này là một trong những người bạn thân thiết của Coco. Một người có vóc dáng nhỏ bé, trông như cần được bảo vệ, nhưng tôi nghe nói sức mạnh của cô ấy cũng không kém cạnh gì - có thể đập vỡ cả một bức tường nếu tức giận. Tôi chỉ hy vọng rằng cái tên này sẽ đỡ hơn chút ít - mà có khi cũng chẳng cần ước đâu, vì tôi đã dự báo trước được mọi chuyện sẽ thế nào mà.

Tôi lật trang hồ sơ, đọc dòng tên đầu tiên. Tôi đã nghe nói về cô ấy - một thiên thần có sức mạnh đáng nể và một trái tim nhân hậu, nhưng cũng có thể rất cứng đầu khi cần. Tôi đặt hồ sơ xuống, nhìn ra cửa sổ, nơi bầu trời tháng Bảy đang dần tối lại.

``\textit{Chào mừng đến với \hll, những cơn bão sắp tới. Tôi chỉ hy vọng là các em sẽ đỡ hơn Coco-chan một chút.}'' tôi thầm nghĩ, vừa cười vừa lắc đầu.

\begin{center}
  {\Large \abv{Bonus}}
\end{center}

\begin{figure}[H]
  \centering
  \includegraphics[width=12.5cm]{../../../../Image/Book?/5-5n.png}
\end{figure}

\end{document}